\begin{table}[!htb]
\caption{Evaluasi Pemenuhan Kebutuhan Nonfungsional}
\label{tab:evaluasi_knf}
\footnotesize
\begin{tabular*}{\textwidth}{@{\extracolsep{\fill}}|>{\raggedright\arraybackslash}m{6.2cm}|>{\centering\arraybackslash}m{1.3cm}|>{\centering\arraybackslash}m{1.5cm}|>{\centering\arraybackslash}m{1.3cm}|>{\centering\arraybackslash}m{1.4cm}|}
\hline
\multicolumn{1}{|>{\centering\arraybackslash}m{6.2cm}|}{\textbf{Kebutuhan}} & \textbf{Cloud} & \textbf{\textit{Open-Source}} & \textbf{Kustom} & \textbf{Hibrida} \\
\hline
KNF-01: Kinerja Latensi Rendah & 2 & 2 & 2 & 2 \\
\hline
KNF-02: \textit{Throughput} Tinggi & 2 & 2 & 2 & 2 \\
\hline
KNF-03: Skalabilitas Horizontal & 2 & 2 & 1 & 2 \\
\hline
KNF-04: Keandalan \& Toleransi Kesalahan & 2 & 2 & 1 & 2 \\
\hline
KNF-05: Konsistensi \& Kebenaran & 1 & 2 & 2 & 1 \\
\hline
KNF-06: \textit{Observability} & 2 & 2 & 1 & 2 \\
\hline
KNF-07: Keamanan & 2 & 2 & 1 & 2 \\
\hline
KNF-08: Ekstensibilitas \& Modularitas & 0 & 2 & 2 & 1 \\
\hline
KNF-09: Akses Konsol Terpadu & 2 & 1 & 0 & 2 \\
\hline
KNF-10: Efisiensi Sumber Daya & 1 & 2 & 1 & 1 \\
\hline
KNF-11: Portabilitas & 0 & 2 & 1 & 1 \\
\hline
KNF-12: Kemudahan Pemeliharaan & 2 & 1 & 0 & 1 \\
\hline
\textbf{Total Skor KNF} & \textbf{18} & \textbf{22} & \textbf{14} & \textbf{19} \\
\hline
\end{tabular*}
\skorlegend
\end{table}